\documentclass[10pt,a4paper]{article}
\usepackage[utf8]{inputenc}
\usepackage[T1]{fontenc}
\usepackage[margin=2.8cm]{geometry}
\usepackage{amsmath,amssymb,amsthm}
\allowdisplaybreaks
\usepackage{mathpartir}
\usepackage{stmaryrd}
\usepackage{listings}
\usepackage{graphicx}
\usepackage{array}
\usepackage{booktabs}
\usepackage[round]{natbib}
\usepackage[hypertexnames=false]{hyperref}
\usepackage[table]{xcolor}
\usepackage{float}
\usepackage{tcolorbox}
\tcbuselibrary{skins,breakable}

% ------------------------------------------------------------------
% Notation macros
% ------------------------------------------------------------------
\newcommand{\desc}{\mathrel{\blacktriangleleft}}      % "correctly described by" relation
\newcommand{\conflict}{\mathrel{\otimes}}             % computed-effect conflict
\newcommand{\s}{\mathrel{\boxtimes}}                  % element-wise conflict
\newcommand{\elemconf}{\s}                              % element-wise conflict (aux relation)
\newcommand{\approxx}{\mathrel{\sqsubseteq}}          % trace covered by computed effect
\newcommand{\sem}[1]{\llbracket #1 \rrbracket}
\newcommand{\Read}{\mathsf{Read}}
\newcommand{\Write}{\mathsf{Write}}
\newcommand{\Alloc}{\mathsf{Alloc}}
\newcommand{\Rs}{\mathsf{R}}
\newcommand{\Ws}{\mathsf{W}}
\newcommand{\As}{\mathsf{A}}
\newcommand{\Rhat}{\widehat{\mathsf{R}}}
\newcommand{\What}{\widehat{\mathsf{W}}}
\newcommand{\Ahat}{\widehat{\mathsf{A}}}
\newcommand{\Readhat}{\widehat{\Read}}
\newcommand{\Writehat}{\widehat{\Write}}
\newcommand{\Top}{\mathsf{Top}}
\newcommand{\Pure}{\mathsf{Pure}}
\newcommand{\Nil}{\mathsf{Nil}}
\newcommand{\tracecat}{\mathbin{\bullet}}
\newcommand{\capp}{\mathsf{capp}\;}
\newcommand{\eapp}{\mathsf{eapp}\;}
\newcommand{\cond}{\mathsf{cond}\;}
\newcommand{\alloc}{\mathsf{alloc}\;}
\newcommand{\getop}{\mathsf{get}\;}
\newcommand{\setop}{\mathsf{set}\;}
\newcommand{\at}{\mathord{@}}
\newcommand{\thetahat}{\widehat{\Theta}}
\newcommand{\ReadOnly}{\mathsf{ReadOnly}}
\newcommand{\Disjoint}{\mathsf{Disjoint}}
\newcommand{\DetTrace}{\mathsf{Det}}
\newcommand{\MapDisjoint}{\mathsf{MapDisjoint}}
\newcommand{\traceset}[1]{\llbracket #1\rrbracket^{\subseteq}}% legacy alias
\newcommand{\dynset}[1]{\lvert #1\rvert_{\mathsf{dyn}}}
\newcommand{\staticset}[1]{\lvert #1\rvert_{\epsilon}}
\newcommand{\cerase}[1]{\lfloor #1 \rfloor_{\epsilon}}
\newcommand{\actleq}{\mathrel{\preccurlyeq_{\Theta}}}
\newcommand{\covers}{\mathrel{\sqsupseteq}}
\newcommand{\silent}{\mathsf{silent}}
\newcommand{\efun}[1]{e^{\mathsf{fun}}_{#1}}
\newcommand{\earg}[1]{e^{\mathsf{arg}}_{#1}}

\newtheorem{theorem}{Theorem}
\newtheorem{lemma}{Lemma}
\newtheorem{corollary}{Corollary}
\newtheorem{definition}{Definition}

\newtcolorbox{formalbox}[1]{
  enhanced jigsaw,
  breakable,
  colback=white,
  colframe=black,
  boxrule=.45pt,
  arc=0pt,
  left=5pt,right=5pt,top=8pt,bottom=5pt,
  title={#1},
  fonttitle=\bfseries,
  coltitle=black,
  attach boxed title to top center={yshift=-2mm},
  boxed title style={colback=white,colframe=white,boxrule=0pt,arc=0pt,left=4pt,right=4pt},
  before skip=1.1em,after skip=1.1em
}


\lstset{basicstyle=\ttfamily\small, keywordstyle=\bfseries,
        morekeywords={fun,let,in,ref,at,region,end}, columns=fullflexible, mathescape=true}

\title{Executable Effect Summaries for Deterministic Parallelism}
\author{Vitor Rodrigues\thanks{VORTEX-Colab, Porto, Portugal.
  \texttt{vitor.rodrigues@vortex-colab.com}}
  \and
  Matthew Fluet\thanks{Rochester Institute of Technology, New York, USA.
  \texttt{mtf@cs.rit.edu}}}
\date{}

\begin{document}
\maketitle

\begin{abstract}
Static effect systems can guarantee deterministic parallelism, but region-level
approximations often reject computations whose concrete run-time accesses are
independent. Dynamic and optimistic systems recover more parallelism by
observing concrete accesses, but typically require logging, validation,
rollback, or replay. This paper explores a middle ground: \emph{surface-effect
terms}, source-level executable summaries that are statically checked against
the computations they describe and evaluated before those computations are
run.

We present a region-polymorphic simply typed $\lambda$-calculus with mutable
references and executable effect summaries. The mechanized development proves
that verified summaries approximate the dynamic traces of their paired
computations, and that terminating continuation-machine executions are
deterministic. Surface effects therefore provide a statically verified,
source-level form of pre-execution access prediction: programmers may expose
value-dependent access information before a computation is run, while the
metatheory ensures that the summary covers the computation trace it claims to
describe.
\end{abstract}

\tableofcontents

\section*{Introduction}

Multicore processors make parallelism increasingly attractive, but shared-state
parallel programs are difficult to reason about. The programmer wants the
performance benefits of running independent computations at the same time, while
the language designer wants the semantic simplicity of deterministic execution:
well-typed programs should not acquire data races, deadlocks, or schedule-
dependent results merely because the implementation exploits parallel hardware.
A central challenge is therefore to expose parallelism without asking the
programmer to reason directly about low-level interleavings.

The language studied in this paper is a region-polymorphic $\lambda$-calculus
with mutable references and executable effect summaries. Typability is established
using a type-and-effect discipline in the Lucassen--Gifford lineage
\citep{lucassen1988}. Region effects and non-interference checks for
deterministic parallelism are not, by themselves, the novelty claimed here; DPJ,
for example, uses modular compile-time type-and-effect checking to guarantee
deterministic semantics \citep{bocchino2009}. The question we ask is more
specific: can the programmer provide a checked, executable summary that is more
precise than a coarse static region effect, but less intrusive than fully
logging or validating every dynamic access?

Throughout the paper, we call the source-level executable effect-summary terms
\emph{surface-effect terms}; when the context is clear, we use \emph{surface
effects} as shorthand for this term language. Evaluating a surface-effect term
yields a \emph{computed effect}, which is the run-time summary value available
to later analyses or runtimes. The contribution is not merely that effects
appear in types, but that a computation can be paired with another well-typed
source term that computes a summary of the computation before the computation
itself is run.

Two degenerate surface-effect terms are useful for intuition. `$\Pure$' denotes
the empty summary and is appropriate for computations that perform no dynamic
memory actions. `$\Top$' is the opposite extreme: it is a sound summary for any
computation, but it gives up precision and therefore prevents a useful
address-level comparison. Between these extremes, surface-effect terms
can compute value-dependent summaries such as concrete reads and writes. This
precision is useful when two computations mention the same static region but
access different concrete locations within that region.

The central intuition is that one well-typed term is described by \emph{another}
well-typed term of the same language. A summary may inspect run-time values and,
for example, compute a concrete region/location pair instead of retaining only a
region-level read or write. Such a summary may distinguish accesses that a
region-only static effect merges. The formal development does not define a
global precision lattice relating static effects, computed effects, and traces,
nor does it prove that a more precise summary is necessarily more expensive. The
soundness claim is instead expressed by explicit approximation relations:
static effects soundly cover dynamic actions, and a verified surface-effect term
evaluates to a computed effect that soundly covers the trace of its paired
computation.

Put simply, the surface-effects approach is a hybrid design point for
value-dependent effect prediction. Purely static approaches are modular and
perform the conflict check at compile time, but they may be conservative when
pointer-rich data structures collapse many concrete locations into the same
abstract region. Purely dynamic and optimistic approaches observe concrete
accesses directly, and hence can be very precise, but they usually require
logging, validation, rollback, or replay. Surface effects place the check in
between: a source-level summary is supplied and verified statically, evaluated
before the paired computation runs, and then used as a more precise dynamic
summary of the computation's future accesses.

Figure~\ref{fig:middleground} previews this design point symbolically.  The
notation is introduced fully in Sections~\ref{sec:effects}--\ref{sec:opsem}.
Here it is only schematic: $s\longrightarrow s';\eta$ denotes one step of the
operational machine together with the trace fragment $\eta$ emitted by that
step, while $s\longrightarrow^{*}s';\phi$ denotes a finite execution whose
fragments have been accumulated into the structured trace $\phi$. Static
checking compares region effects $\epsilon_i$. Surface-effect checking first
verifies source expressions against surface-effect terms $\thetahat_i$, then
evaluates those terms by the same operational machine to computed effects
$\Theta_i$. Dynamic approaches compare traces $\phi_i$ obtained from the
actual executions. The table writes $\otimes$ for the relevant comparison at
each level; the formal development later refines the computed-effect comparison
into computed-effect disjointness and non-conflict premises.

Static effects are inexpensive and available at compile time, but their
region-level approximation may reject independent concrete accesses. Surface
effects pay the cost of evaluating executable summaries before running the
paired computations, but can recover value-dependent precision. Dynamic tracking
observes the concrete accesses themselves, but usually requires logging,
validation, rollback, or replay.


\begin{figure}[t]
\centering
\scriptsize
\setlength{\tabcolsep}{4pt}
\renewcommand{\arraystretch}{1.13}
\resizebox{\textwidth}{!}{%
\begin{tabular}{@{}>{\raggedright\arraybackslash}p{.20\textwidth}
                p{.29\textwidth}c p{.29\textwidth}@{}}
\hline
\multicolumn{4}{c}{$e_1 \quad e_2$}\\[-.2ex]
\multicolumn{4}{c}{\bfseries Three places to compare effects}\\[.4ex]
\hline
\textbf{Static region effects}
  & $\Gamma \vdash e_1 : \tau_1\,!\,\epsilon_1$
  & $\epsilon_1 \conflict \epsilon_2$
  & $\Gamma \vdash e_2 : \tau_2\,!\,\epsilon_2$\\[.5ex]
\rowcolor{gray!12}
\textbf{Surface effects}
  & $\Gamma;\Delta;P \vdash e_1 \desc \thetahat_1$
  & $\Theta_1 \conflict \Theta_2$
  & $\Gamma;\Delta;P \vdash e_2 \desc \thetahat_2$\\
\rowcolor{gray!12}
  & $\thetahat_1 \longrightarrow^{*} \Theta_1$
  &
  & $\thetahat_2 \longrightarrow^{*} \Theta_2$\\
\rowcolor{gray!12}
  & $\phi_1 \approxx \Theta_1$
  &
  & $\phi_2 \approxx \Theta_2$\\[.5ex]
\textbf{Dynamic traces}
  & $e_1 \leadsto \phi_1$
  & $\phi_1 \conflict \phi_2$
  & $e_2 \leadsto \phi_2$\\
\hline
\end{tabular}%
}
\caption{Surface effects as a middle ground between static region effects and
  dynamic access tracking.}
\label{fig:middleground}
\end{figure}

Operationally, the ordinary continuation-machine relation
$s\longrightarrow s';\eta$ evaluates expression structure and emits the trace
fragment $\eta$. A surface-effect term is evaluated by this same machine; there
is no separate summary evaluator. Thus the paper's runtime theorem is an
operational theorem about finite small-step executions, not merely a statement
about typing rules.

The static semantics uses type-and-effect judgments of the form
$\Gamma \vdash e : \tau\,!\,\epsilon$, where $\Gamma$ is a typing context,
$e$ is an expression, $\tau$ is a type, and $\epsilon$ is a static effect (see
Section~\ref{sec:staticsem}).  The static semantics of surface effects defines
the back-triangle relation ($\desc$), which states that a computational
expression $e$ is correctly described by a surface-effect term $\thetahat$.
At run time, that term is evaluated by the same small-step machine as every
other expression. When it returns an effect value, that value is the computed
effect $\Theta$. If a computed effect is too imprecise, for example because a
summary evaluated to $\top$, the result remains sound but conservative.

Static region/effect systems such as DPJ make non-interference a modular
compile-time obligation, but region-level summaries may merge concrete accesses
that are distinct at run time \citep{bocchino2009}. Runtime inspector--executor
techniques address irregular dependences by discovering conflict-free work
before executing it in parallel \citep{arenaz2004}. Surface effects occupy a
related middle position: the access summary is supplied as a typed source-
language term, checked against the computation, and evaluated before the
described computation is run. Unlike optimistic post-validation, the calculus
does not start the described computation and then roll it back after discovering
a summary conflict; however, this fact alone does not establish that
pre-checking is cheaper in total.

The formal contribution is therefore a sound scheduling interface rather than
an asymptotic performance result. At compile time, an inductive correctness
relation validates programmer-written computation/summary pairs. At run time,
the summary terms evaluate to computed effects that may contain concrete
region/location information. The mechanized development connects these stages:
summary correctness gives trace approximation, summary-level safety gives a
deterministic parallel trace, and deterministic trace replay gives equivalent
final heaps. False positives remain possible because summaries are
approximations and `$\Top$' is deliberately maximally conservative.

The contributions of this paper are:

\begin{itemize}
  \item an executable effect-summary mini-language embedded in a
        region-polymorphic calculus with mutable references;
  \item an inductive correctness relation for statically verifying a
        computation against a programmer-written, value-dependent summary term;
  \item a trace-emitting small-step operational semantics in which executable
        summaries are ordinary expressions of type $\mathsf{eff}$;
  \item a mechanized approximation theorem showing that a verified summary
        soundly covers the dynamic trace of its computation;
  \item a mechanized continuation-machine safety layer showing that every
        finite prefix of a well-typed execution remains type safe and not
        stuck, plus explicit terminal determinism and terminal correctness
        theorems for ordinary executions.
\end{itemize}

The rest of this paper is organized as follows. We start by presenting a
motivating example in Section~\ref{sec:motivating}.
Section~\ref{sec:effects} distinguishes static effects, dynamic traces,
surface-effect terms, and computed effects. The syntax for types and the
surface language, including the surface-effect language, is given in
Section~\ref{sec:syntax}. The static semantics of the type-and-effect system is
presented in Section~\ref{sec:staticsem}, and the operational semantics of the
ordinary continuation machine is presented in Section~\ref{sec:opsem}. The
computed-effect conflict and disjointness are introduced in Section~\ref{sec:effects}
before presenting the correctness theorem for surface-effect terms in
Section~\ref{sec:correctness}. The small-step finite-prefix safety layer is
summarized in Section~\ref{sec:smallstep}. We conclude in
Section~\ref{sec:conclusions} after discussion of related work.
\section{Motivating Example}
\label{sec:motivating}

The ML-like program in Fig.~\ref{fig:example} illustrates the intuition behind
surface-effect terms. The function \lstinline|incr| increments the integer
stored in a mutable reference. Its argument has type
$\mathsf{ref}\,\mathsf{int}\;\mathsf{at}\;\rho$, and the latent static effect of
the body is $\{\Rs(\rho),\Ws(\rho)\}$: the reference is both read and written.
A conventional region-effect check therefore treats two calls to \lstinline|incr|
at the same region as potentially conflicting.

This is conservative. A list of references is usually homogeneous in its region
parameter, so two variables \lstinline|x| and \lstinline|y| drawn from the same
list may both have type $\mathsf{ref}\,\mathsf{int}\;\mathsf{at}\;\rho$. A
static region effect cannot distinguish whether \lstinline|x| and \lstinline|y|
are aliases or two different addresses inside the same region. Surface effects
are intended to recover that concrete-address information before the described
computation is run.

\begin{figure}[ht]
\begin{lstlisting}
fun incr [rho] (r : int ref at rho) : unit =
  r := !r + 1

fun incr_eff [rho] (r : int ref at rho) : eff =
  $\widehat{\Read(r)} \oplus \widehat{\Write(r)}$

fun use (x : int ref at 7) : unit =
  incr[7] x       (* described by incr_eff[7] x *)
\end{lstlisting}
\caption{Motivating example: a computation paired with an executable
surface-effect summary.}
\label{fig:example}
\end{figure}

The role of \lstinline|incr_eff| is to compute the summary that the static
region effect cannot express. The surface-effect term
$\widehat{\Read(r)} \oplus \widehat{\Write(r)}$ is evaluated as an ordinary
effect-valued expression. Its result is a
computed effect that mentions the concrete region/location pair of the actual
argument. Thus, two calls whose arguments are distinct locations in the same
region can produce non-conflicting computed effects, while two calls on the same
location still conflict.

The point of the example is not that every invocation of \lstinline|incr| is
safe to parallelize. Statically, each call has the region effect
$\{\Rs(\rho),\Ws(\rho)\}$, so two calls at the same region appear to conflict.
The surface-effect term \lstinline|incr_eff| refines that coarse description by
computing the summary $\widehat{\Read(r)} \oplus \widehat{\Write(r)}$ from the actual
argument \lstinline|r|. The computed effect can mention concrete
region/location pairs rather than only the region $\rho$.

The same example also illustrates the precision spectrum. A pure computation
may be summarized by `$\Pure$', the empty surface-effect term. The computation
\lstinline|incr| cannot soundly use `$\Pure$', because it reads and writes its
argument. It could, however, always use `$\Top$': this summary is sound but
maximally conservative, so any later consumer of the summary would learn no
useful address-level information.
The useful summary lies between these extremes: $\widehat{\Read(r)} \oplus
\widehat{\Write(r)}$ is still sound, but it exposes enough run-time address information
to avoid some false conflicts.

Informally, if \lstinline|incr_eff| correctly describes \lstinline|incr|, then
the static relation $(\mathsf{incr} \desc \mathsf{incr\_eff})$ must hold. The
precision slogan is
\[
  \epsilon \covers \Theta \covers \phi,
\]
but this is not a single global order. It abbreviates two typed bridges: a
static effect $\epsilon$ covers the region erasure of a computed effect
$\Theta$, and the computed effect $\Theta$ covers the dynamic trace $\phi$ by
the relation $\phi\approxx\Theta$. In the motivating case, the
programmer-written summary is checked once as part of the program, but its
value-dependent result is computed at run time by the ordinary machine.

Finally, surface-effect terms used as summaries are required to be read-only.
The purpose of \lstinline|incr_eff| is to inspect the reference value well enough
to construct a summary, not to perform the update that \lstinline|incr| will
perform. This restriction is what lets the metatheory separate summary
evaluation from the computation it describes: evaluating the summary can affect
the scheduling decision, but it cannot itself perform the writes whose safety is
being checked.
\section{Static Effects, Dynamic Traces, Surface Effects, and Computed Effects}
\label{sec:effects}

This section fixes the vocabulary that is only previewed in the Introduction.
The development uses four related, but deliberately distinct, notions of
``effect'': static effects, dynamic traces, surface-effect terms, and computed
effects. The distinction is essential. Static effects are compile-time
approximations; dynamic traces are observations of execution; surface-effect
terms are source syntax; and computed effects are the values obtained by running
that syntax.

\paragraph{Regions, addresses, and actions.}
We write $\rho$ for a region variable, $r$ for a run-time region constant,
$l$ for a heap location, and $l\at r$ for the concrete address consisting of
location $l$ in region $r$. A static region $p$ may be either a region variable
or a region constant. Static actions describe allocation, read, and write at the
level of regions, whereas dynamic actions describe concrete events at allocated
addresses. The value component of a dynamic action is important for replaying a
trace, but is usually irrelevant for conflict checking, so examples often omit
it.

\begin{formalbox}{Effect Definitions}
\small
\[
\begin{array}{rcll}
\rho &\in& \mathit{RgnVar} & \text{region variables}\\
r &\in& \mathit{RgnConst} & \text{region constants}\\
l &\in& \mathit{Loc} & \text{heap locations}\\
l\at r &\in& \mathit{Addr} \triangleq \mathit{Loc}\times\mathit{RgnConst}
  & \text{concrete heap addresses}\\[.4ex]
p &::=& \rho \mid r & \text{static regions}\\
\alpha &::=& \As(p) \mid \Rs(p) \mid \Ws(p) & \text{static actions}\\
\epsilon &\in& \mathcal{P}(\alpha) & \text{static effects}\\[.4ex]
\iota &::=& \Alloc(l\at r,v) \mid \Read(l\at r,v) \mid \Write(l\at r,v)
  & \text{dynamic actions}\\
\phi &::=& \Nil \mid \iota \mid \phi_1\tracecat\phi_2 \mid \phi_1\parallel\phi_2
  & \text{dynamic traces}\\[.4ex]
\theta &::=& \Ahat(r) \mid \Rhat(r) \mid \What(r)
       \mid \widehat{\Read(l\at r)} \mid \widehat{\Write(l\at r)}
  & \text{computed actions}\\
\Theta &::=& \top \mid \{\theta_1,\ldots,\theta_n\}
  & \text{computed effects}
\end{array}
\]
\end{formalbox}

\paragraph{Static effects and dynamic traces.}
A static effect $\epsilon$ is the effect component of a type-and-effect judgment
$\Gamma \vdash e : \tau\,!\,\epsilon$. It abstracts over all possible executions
of $e$ using region-level actions such as $\Rs(\rho)$ or $\Ws(\rho)$. A dynamic
trace $\phi$ is produced by the operational semantics. It records the concrete
allocation, read, and write events emitted by a particular execution. Structured
traces compose sequentially and in parallel; this structure is used by the
trace-replay and deterministic-parallelism theorems.

\paragraph{Surface-effect terms and computed effects.}
A surface-effect term $\thetahat$ is an expression of the source language whose
result type is $\mathsf{eff}$. Evaluating such a term produces a computed effect
$\Theta$. Computed effects intentionally mix abstract region actions and
concrete address actions: a summary may say only that a computation writes some
location in region $r$, or it may say that the computation writes the particular
address $l\at r$. This is the mechanism that lets a summary be more precise than
a static region effect while still being evaluated before the paired computation
runs.

Not every static or dynamic action has the same status as a computed action.
Concrete reads and writes can appear in computed effects, because a summary term
can evaluate an argument reference and obtain its region/location pair. Concrete
allocation is not used as a computed action in the mechanization; allocation is
summarized abstractly at the region level. Conversely, dynamic traces contain
concrete allocations, reads, and writes because they describe what an execution
actually did.

The surface-effect term `$\Top$' evaluates to the distinguished computed effect
$\top$. It is a sound but maximally conservative approximation: every dynamic
trace is covered by $\top$, but $\top$ conflicts with every computed effect
and cannot satisfy the computed-effect disjointness premise. Thus, a
surface-effect term may always fall back to `$\Top$' for soundness, but doing so
loses all useful disjointness information. By contrast, `$\Pure$'
evaluates to the empty computed-effect set and describes computations with no
dynamic actions.

The useful slogan from the earlier presentation is
$\epsilon\covers\Theta\covers\phi$: static effects are coarser than computed
effects, and computed effects are coarser than the traces they summarize.  In
this paper the slogan is only a guide to the reader; formally, the objects live
in different syntactic categories and are related by typed bridges rather than
by a single global lattice.


\begin{formalbox}{Typed Approximation Bridges}
\small
The slogan $\epsilon\covers\Theta\covers\phi$ abbreviates two typed
approximations.  The left bridge compares a static effect with the
\emph{static erasure} of a finite computed effect:
\[
\begin{array}{rclcrcl}
\cerase{\Ahat(r)} &=& \As(r)
&\qquad&
\cerase{\Rhat(r)} &=& \Rs(r)\\[0.5mm]
\cerase{\What(r)} &=& \Ws(r)
&\qquad&
\cerase{\widehat{\Read(l\at r)}} &=& \Rs(r)\\[0.5mm]
\cerase{\widehat{\Write(l\at r)}} &=& \Ws(r).
\end{array}
\]
\[
\cerase{\{\theta_1,\ldots,\theta_n\}}
  = \{\cerase{\theta_1},\ldots,\cerase{\theta_n}\},
\qquad
\epsilon\covers\Theta
\quad\text{iff}\quad
\Theta\neq\top\text{ and }\cerase{\Theta}\subseteq\epsilon.
\]
The right bridge compares a computed effect with a structured dynamic trace.
It first forgets only the intensional trace structure:
\[
\begin{array}{rclcrcl}
\dynset{\Nil} &=& \emptyset
&\qquad&
\dynset{\iota} &=& \{\iota\}\\[0.5mm]
\dynset{\phi_1\tracecat\phi_2} &=&
  \dynset{\phi_1}\cup\dynset{\phi_2}
&\qquad&
\dynset{\phi_1\parallel\phi_2} &=&
  \dynset{\phi_1}\cup\dynset{\phi_2}.
\end{array}
\]
This dynamic-action view still contains concrete locations and values, so it is
related to computed effects by the structural approximation judgment
$\phi\approxx\Theta$, not by subset inclusion.  If we additionally erase
dynamic actions to ordinary static actions,
\[
\begin{array}{rclcrcl}
\cerase{\Alloc(l\at r,v)} &=& \As(r)
&\qquad&
\cerase{\Read(l\at r,v)} &=& \Rs(r)\\[0.5mm]
\cerase{\Write(l\at r,v)} &=& \Ws(r)
&\qquad&
\staticset{\phi} &=& \{\cerase{\iota}\mid \iota\in\dynset{\phi}\},
\end{array}
\]
then subset inclusion becomes well typed at the static level:
\[
\phi\approxx\Theta\quad\text{and}\quad\Theta\neq\top
\quad\Longrightarrow\quad
\staticset{\phi}\subseteq\cerase{\Theta}.
\]
Consequently, when both bridges hold for a finite computed effect, the informal
three-level comparison can be read as
\[
\staticset{\phi}\subseteq\cerase{\Theta}\subseteq\epsilon.
\]
\end{formalbox}

\paragraph{Computed-effect coverage.}
The right-hand bridge, $\Theta\covers\phi$, is made precise by the
structural trace-to-computed-effect judgment $\phi\approxx\Theta$.  This
judgment does not evaluate a trace and does not turn an execution trace into a
computed effect.  Instead, it checks whether a computed effect covers each
dynamic action that can occur in the structured trace.  The action case
delegates to an action-in-summary relation $\iota\actleq\Theta$.

\begin{formalbox}{Action Coverage}
\footnotesize
\noindent\textbf{Action-in-summary coverage}\hfill $\boxed{\iota\actleq\Theta}$
\begin{mathpar}
\inferrule*[lab=Act-Top]{ }
  {\iota \actleq \top}

\inferrule*[lab=Act-Alloc-Abs]
  {\Ahat(r) \in \vartheta}
  {\Alloc(l\at r,v) \actleq \vartheta}

\inferrule*[lab=Act-Read-Conc]
  {\widehat{\Read(l\at r)} \in \vartheta}
  {\Read(l\at r,v) \actleq \vartheta}

\inferrule*[lab=Act-Read-Abs]
  {\Rhat(r) \in \vartheta}
  {\Read(l\at r,v) \actleq \vartheta}

\inferrule*[lab=Act-Write-Conc]
  {\widehat{\Write(l\at r)} \in \vartheta}
  {\Write(l\at r,v) \actleq \vartheta}

\inferrule*[lab=Act-Write-Abs]
  {\What(r) \in \vartheta}
  {\Write(l\at r,v) \actleq \vartheta}
\end{mathpar}
\end{formalbox}

\begin{formalbox}{Trace Coverage}
\footnotesize
\noindent\textbf{Structural trace-to-computed-effect coverage}\hfill $\boxed{\phi\approxx\Theta}$
\begin{mathpar}
\inferrule*[lab=Trace-Nil]{ }
  {\Nil \approxx \Theta}

\inferrule*[lab=Trace-Elem]
  {\iota \actleq \Theta}
  {\iota \approxx \Theta}

\inferrule*[lab=Trace-Seq]
  {\phi_1\approxx\Theta \\ \phi_2\approxx\Theta}
  {\phi_1\tracecat\phi_2 \approxx \Theta}

\inferrule*[lab=Trace-Par]
  {\phi_1\approxx\Theta \\ \phi_2\approxx\Theta}
  {\phi_1\parallel\phi_2 \approxx \Theta}
\end{mathpar}
\end{formalbox}

For finite computed effects, the structural judgment has the expected
action-set characterization:
\[
\phi\approxx\Theta
\quad\Longleftrightarrow\quad
\forall\iota\in\dynset{\phi}.\;\iota\actleq\Theta.
\]
This equivalence explains the role of the two trace views above.  The dynamic
view $\dynset{\phi}$ keeps concrete locations and values and therefore compares
with computed effects through $\approxx$, not through raw subset inclusion.
Only after static erasure do ordinary subset comparisons become well typed: if
$\phi\approxx\Theta$ and $\Theta\neq\top$, then
$\staticset{\phi}\subseteq\cerase{\Theta}$.

The case $\Theta=\top$ is the conservative escape hatch for this right-hand
bridge: it covers every dynamic trace, but it is not a finite computed-action
set whose erasure recovers useful static precision.  The mechanization defines
$\top$, the empty computed effect, and union of computed effects, but it does
not require a single lattice over $\epsilon$, $\Theta$, and $\phi$.

The left-hand bridge $\epsilon\covers\Theta$ therefore uses the erasure morphism
$\cerase{\cdot}$ to forget the extra precision of computed effects. It erases
both the abstract/concrete distinction and concrete locations: for example,
$\widehat{\Read(l\at r)}$ and $\Rhat(r)$ both collapse to the static region
action $\Rs(r)$.

The right-hand bridge keeps the concrete actions in the trace and is expressed
by $\phi\approxx\Theta$. This is the bridge established by surface-effect
correctness: evaluating a verified surface-effect term produces a computed
effect that covers the trace of the paired computation. A later client can
combine this bridge with computed-effect disjointness and non-conflict when it
wants to reason about separation of two traces.

\paragraph{Computed-effect conflict and disjointness.}\label{sec:disjoint}
The computed-effect guard is stated entirely over computed effects. It requires the
universal disjointness predicate $\Disjoint(\Theta_1,\Theta_2)$ and the absence
of the existential conflict predicate $\Theta_1\conflict\Theta_2$.  We write
$\theta_1\s\theta_2$ for element-wise conflict between computed actions, using
the square-cross symbol $\s$, and $\Theta_1\conflict\Theta_2$ for conflict
between computed effects.

\begin{formalbox}{Existence of conflicts ($\s$) in $\theta$}
\footnotesize
\[
\boxed{\theta_1\s\theta_2}
\]
\begin{mathpar}
\inferrule*[lab=C-READ-WRITE]
  {\theta = \Readhat(l\at r) \,\vee\, \theta = \Writehat(l\at r) \,\vee\,
   \theta = \Rhat(r) \,\vee\, \theta = \What(r)}
  {\What(r)\s\theta \;\wedge\; \Writehat(l\at r)\s\theta}
\end{mathpar}
\end{formalbox}

\begin{formalbox}{Existence of conflicts ($\otimes$) in $\Theta$}
\footnotesize
\[
\boxed{\Theta_1\conflict\Theta_2}
\]
\begin{mathpar}
\inferrule*[lab=C-EXT-THETA]
  {\theta_1 \in \Theta_1 \\ \theta_2 \in \Theta_2 \\
   \theta_1 \s \theta_2}
  {\Theta_1 \conflict \Theta_2}

\inferrule*[lab=C-TOP]{ }
  {\Theta \conflict \top}
\end{mathpar}
\end{formalbox}

\section{Syntax}
\label{sec:syntax}

The surface language is a region-polymorphic, call-by-value simply typed
$\lambda$-calculus with mutable references and recursive functions. The
distinctive syntactic form is the recursive abstraction
$\mu f.\lambda x.\{e_b\mid e_s\}$, which packages two source expressions: the
computational body $e_b$ and the surface-effect body $e_s$.  Computational
application, written $\capp e_f\,e_a$, runs $e_b$; effect application, written
$\eapp e_f\,e_a$, runs $e_s$ and must produce a value of type
$\mathsf{eff}$.

The grammar below presents the surface language in a compact, framed form. It uses the
current terminology.  A surface-effect term is not a separate run-time language:
it is source syntax of type $\mathsf{eff}$.  The meta-variable $\thetahat$ only
isolates the expressions normally used as executable summaries.  In the
mechanized syntax, $\Pure$ is represented by the empty effect expression, while
$\Top$ is the conservative effect expression.

\begin{formalbox}{Surface Language Syntax}
\small
\begin{align*}
\textit{Regions}\quad
p ::= {}& \rho \mid r.\\[.7ex]
\textit{Types}\quad
\tau ::= {}& \mathsf{nat}
  \mid \mathsf{bool}
  \mid \mathsf{unit}
  \mid \mathsf{eff}\\
& \mid \mathsf{ref}\;p\;\tau
  \mid \tau_x \to \{\tau_c\,!\,\epsilon_c \mid
                     \mathsf{eff}\,!\,\epsilon_s\}
  \mid \forall \rho.\tau\,!\,\epsilon.\\[.7ex]
\textit{Expressions}\quad
 e ::= {}& x \mid n \mid b
  \mid \mu f.\lambda x.\{e_b\mid e_s\}
  \mid \Lambda \rho.e\\
& \mid \capp e\,e
  \mid \eapp e\,e
  \mid e[p]
  \mid \cond(e,e,e)\\
& \mid \alloc p\,e
  \mid \getop p\,e
  \mid \setop p\,e\,e\\
& \mid e+e \mid e-e \mid e*e \mid e=e\\
& \mid \thetahat.\\[.7ex]
\textit{Surface-effect terms}\quad
\thetahat ::= {}& \Pure
  \mid \Ahat(p)
  \mid \Rhat(p)
  \mid \What(p)\\
& \mid \widehat{\Read(e)}
  \mid \widehat{\Write(e)}
  \mid \thetahat_1\oplus\thetahat_2
  \mid \Top.
\end{align*}
\end{formalbox}

Types follow the region-polymorphic simply typed calculus used throughout the
formalization, with two additions made explicit here. The base types $\mathsf{nat}$,
$\mathsf{bool}$, and $\mathsf{unit}$ classify numbers, booleans, and the result
of destructive updates; $\mathsf{eff}$ classifies computed-effect values
$\Theta$. A reference type $\mathsf{ref}\;p\;\tau$ describes a mutable cell of
contents type $\tau$ allocated in region $p$, where $p$ may be a region variable
or a concrete region. Function types record two latent effects: in
$\tau_x \to \{\tau_c\,!\,\epsilon_c \mid \mathsf{eff}\,!\,\epsilon_s\}$,
$\epsilon_c$ is the static effect of the computational body and $\epsilon_s$ is
the static effect of the executable summary body. Region-polymorphic types
$\forall \rho.\tau\,!\,\epsilon$ abstract over regions and record the latent
effect of the body under that abstraction.

Expressions include variables, numeric and Boolean constants, recursive
computation/summary abstractions, region abstractions, computational
applications, surface-effect applications, region applications, conditionals,
arithmetic and equality, heap operations, and surface-effect terms. In
$\mu f.\lambda x.\{e_b\mid e_s\}$, the argument
variable $x$ and recursive variable $f$ are available in both $e_b$ and $e_s$.
The type system later requires the summary body to be read-only when it is used
as the executable summary for a computation.

Returning to Fig.~\ref{fig:example}, the function \lstinline|incr| has a
computation body that reads and writes its argument, while \lstinline|incr_eff|
is the corresponding surface-effect body.  Informally, the static correctness
relation must establish
$\Gamma;\Delta;P\vdash \mathsf{incr}\desc \mathsf{incr\_eff}$.
The type of the computation records the coarse static region effect
$\{\Rs(\rho),\Ws(\rho)\}$, whereas evaluating the surface-effect body on a
particular argument can produce a computed effect mentioning the concrete
region/location pair of that argument.  The next section gives the small-step
semantics that performs this evaluation.



\section{Operational Semantics}
\label{sec:opsem}

The operational semantics has one primitive execution relation: the
continuation-machine step $s\longrightarrow s'\,;\,\eta$.  The other judgments
below are derived or proof-facing relations built around that step.  They are
kept distinct because they play different roles in the metatheory.

\begin{center}
\small
\begin{tabular}{@{}>{\raggedright\arraybackslash}p{.18\linewidth}
                  >{\raggedright\arraybackslash}p{.33\linewidth}
                  >{\raggedright\arraybackslash}p{.39\linewidth}@{}}
\toprule
\textbf{Layer} & \textbf{Judgment} & \textbf{Role}\\
\midrule
machine step
  & $s\longrightarrow s'\,;\,\eta$
  & one continuation-machine step emitting one trace fragment\\
finite execution
  & $s\longrightarrow^{*}s'\,;\,\phi$
  & accumulated machine trace\\
trace replay
  & $(\phi,H)\Rightarrow^{*}(\Nil,H')$
  & replay of a produced trace over a heap\\
\bottomrule
\end{tabular}
\end{center}

The notation is intentionally separated.  Machine execution \emph{produces}
trace fragments; the finite-run relation concatenates them into a structured
trace; replay later \emph{consumes} an already produced trace.  A one-step
transition emits only a fragment $\eta$: administrative steps emit $\Nil$, while
heap operations emit singleton fragments such as $\Alloc(l\at r,v)$,
$\Read(l\at r,v)$, and $\Write(l\at r,v)$.  The trace of a longer execution is
not stored in the state or continuation.

\begin{formalbox}{Operational Domains}
\small
\begin{tabular}{@{}>{$}l<{$} @{\quad} l@{}}
E\in\mathit{Env}\triangleq \mathit{Variables}\rightharpoonup\mathit{Values} & (variable assignment)\\
P\in\mathit{Rgn}\triangleq \mathit{RgnVars}\rightharpoonup\mathit{RgnConst} & (regions in scope)\\
H\in\mathit{Heap}\triangleq \mathit{Addr}\rightharpoonup\mathit{Values} & (heap assignment)\\
\langle E,P,e\rangle\in\mathit{Closures}\triangleq
  \mathit{Env}\times\mathit{Rgn}\times\mathit{Expressions} & (function closures)\\
v\in\mathit{Values}\triangleq
  \mathbb{N}+\mathbb{B}+\mathit{Addr}+\mathit{Closures}+\Theta+()+ (v,v) & (dynamic values)\\[.6ex]
s ::= \langle H,E,P,e,k\rangle\mid\langle H,v,k\rangle\mid\mathsf{done}(H,v) & (machine states)\\
a ::= \Alloc(l\at r,v)\mid\Read(l\at r,v)\mid\Write(l\at r,v) & (dynamic actions)\\
\eta ::= \Nil\mid a & (one-step trace fragments)\\
\phi ::= \Nil\mid a\mid \phi\tracecat\phi\mid \phi\parallel\phi & (structured traces)
\end{tabular}
\end{formalbox}

\begin{formalbox}{Finite Machine Executions}
\footnotesize
\noindent\textbf{Finite execution}\hfill $\boxed{s\longrightarrow^{*}s'\,;\,\phi}$

Each finite execution accumulates one-step fragments by induction.
\begin{mathpar}
\inferrule*[lab=TR-Refl]{ }
  {s\longrightarrow^{*}s\,;\,\Nil}

\inferrule*[lab=TR-Step]
  {s\longrightarrow s_1\,;\,\eta \\
   s_1\longrightarrow^{*}s_2\,;\,\phi}
  {s\longrightarrow^{*}s_2\,;\,\eta\tracecat\phi}
\end{mathpar}
Here $\tracecat$ is structured trace sequencing and $\Nil$ is its identity.
Ordinary expression evaluation is the special case
\[
\mathsf{init}(H,E,P,e)
  \longrightarrow^{*}
\mathsf{done}(H',v)\,;\,\phi.
\]
A surface-effect term is just an ordinary expression of type $\mathsf{eff}$;
there is no separate summary evaluator.
\end{formalbox}

\begin{formalbox}{Trace Replay}
\footnotesize
\noindent\textbf{Replay}\hfill $\boxed{(\phi,H)\Rightarrow^{*}(\Nil,H')}$

Replay interprets an already produced structured trace as a heap transformer.
The following selected rules fix the intended direction: execution produces
traces, while replay consumes them.
\begin{mathpar}
\inferrule*[lab=R-Nil]{ }
  {(\Nil,H)\Rightarrow^{*}(\Nil,H)}

\inferrule*[lab=R-Alloc]{ }
  {(\Alloc(l\at r,v),H)\Rightarrow^{*}(\Nil,H[(r,l)\mapsto v])}

\inferrule*[lab=R-Read]{H(r,l)=v}
  {(\Read(l\at r,v),H)\Rightarrow^{*}(\Nil,H)}

\inferrule*[lab=R-Write]{H(r,l)\neq\bot}
  {(\Write(l\at r,v),H)\Rightarrow^{*}(\Nil,H[(r,l)\mapsto v])}

\inferrule*[lab=R-Seq]
  {(\phi_1,H)\Rightarrow^{*}(\Nil,H_1) \\
   (\phi_2,H_1)\Rightarrow^{*}(\Nil,H_2)}
  {(\phi_1\tracecat\phi_2,H)\Rightarrow^{*}(\Nil,H_2)}
\end{mathpar}
The replay relation for parallel traces allows the two subtraces to be consumed
by an interleaving of their replay steps; the determinism theorem later shows
that, under the computed-effect guard, terminating replays agree on the final
heap.
\end{formalbox}


The initial state for expression $e$ is
$\mathsf{init}(H,E,P,e)=\langle H,E,P,e,\mathsf{doneK}\rangle$; terminal
states $\mathsf{done}(H,v)$ have no outgoing machine step.  Continuations are
bookkeeping devices for the call-by-value evaluation order.
A frame records the expression still to be evaluated, the environment and region
map needed to resume evaluation, and the continuation to resume afterward.
Application frames first remember the unevaluated argument while the function
position is evaluated, then remember the closure while the argument is
evaluated, and finally enter either the computational body or the summary body
with the recursive name and argument bound.  Heap-operation frames remember the
pending operation and the region environment while a location or value
subexpression is evaluated.  The trace emitted while evaluating functions,
arguments, or summaries is the machine trace of that evaluation.  A computed
effect is the value returned by evaluating the corresponding summary expression.

\begin{formalbox}{Small-Step Operational Semantics}
\footnotesize
\noindent\textbf{One machine step}\hfill $\boxed{s\longrightarrow s'\,;\,\eta}$
\begin{mathpar}
\inferrule*[lab=S-Const]{ }
  {\langle H,E,P,\mathsf{const}\;n,k\rangle
   \longrightarrow
   \langle H,n,k\rangle\,;\,\Nil}

\inferrule*[lab=S-Var]{E(x)=v}
  {\langle H,E,P,x,k\rangle
   \longrightarrow
   \langle H,v,k\rangle\,;\,\Nil}

\inferrule*[lab=S-Done]{ }
  {\langle H,v,\mathsf{doneK}\rangle
   \longrightarrow
   \mathsf{done}(H,v)\,;\,\Nil}

\inferrule*[lab=S-Clos]
  {\varphi=\mu f.\lambda x.\{e_b\mid e_s\}\;\vee\;\varphi=\Lambda\rho.e}
  {\langle H,E,P,\varphi,k\rangle
   \longrightarrow
   \langle H,\langle E,P,\varphi\rangle,k\rangle\,;\,\Nil}

\inferrule*[lab=S-Capp-Fun]{ }
  {\langle H,E,P,\capp e_f\,e_a,k\rangle
   \longrightarrow
   \langle H,E,P,e_f,K^\mu_f(e_a,E,P,k)\rangle\,;\,\Nil}

\inferrule*[lab=S-Capp-Arg]
  {v_f=\langle E',P',\mu f.\lambda x.\{e_b\mid e_s\}\rangle}
  {\langle H,v_f,K^\mu_f(e_a,E,P,k)\rangle
   \longrightarrow
   \langle H,E,P,e_a,K^\mu_a(E',P',f,x,e_b,e_s,k)\rangle\,;\,\Nil}

\inferrule*[lab=S-Capp-Body]
  {E''=E'[f\mapsto v_f][x\mapsto v_a] \\
   v_f=\langle E',P',\mu f.\lambda x.\{e_b\mid e_s\}\rangle}
  {\langle H,v_a,K^\mu_a(E',P',f,x,e_b,e_s,k)\rangle
   \longrightarrow
   \langle H,E'',P',e_b,k\rangle\,;\,\Nil}

\inferrule*[lab=S-Eapp-Fun]{ }
  {\langle H,E,P,\eapp e_f\,e_a,k\rangle
   \longrightarrow
   \langle H,E,P,e_f,K^e_f(e_a,E,P,k)\rangle\,;\,\Nil}

\inferrule*[lab=S-Eapp-Arg]
  {v_f=\langle E',P',\mu f.\lambda x.\{e_b\mid e_s\}\rangle}
  {\langle H,v_f,K^e_f(e_a,E,P,k)\rangle
   \longrightarrow
   \langle H,E,P,e_a,K^e_a(E',P',f,x,e_b,e_s,k)\rangle\,;\,\Nil}

\inferrule*[lab=S-Eapp-Body]
  {E''=E'[f\mapsto v_f][x\mapsto v_a] \\
   v_f=\langle E',P',\mu f.\lambda x.\{e_b\mid e_s\}\rangle}
  {\langle H,v_a,K^e_a(E',P',f,x,e_b,e_s,k)\rangle
   \longrightarrow
   \langle H,E'',P',e_s,k\rangle\,;\,\Nil}

\inferrule*[lab=S-Alloc]
  {P(p)=r \qquad \mathsf{fresh}(H,r)=l}
  {\langle H,v,K^{ref}(p,P,k)\rangle
   \longrightarrow
   \langle H[(r,l)\mapsto v],l\at r,k\rangle\,;\,\Alloc(l\at r,v)}

\inferrule*[lab=S-Get]
  {P(p)=r \qquad H(r,l)=v}
  {\langle H,l\at r,K^{get}(p,P,k)\rangle
   \longrightarrow
   \langle H,v,k\rangle\,;\,\Read(l\at r,v)}

\inferrule*[lab=S-Set]
  {P(p)=r \qquad H(r,l)\neq\bot}
  {\langle H,v,K^{set}(p,l,P,k)\rangle
   \longrightarrow
   \langle H[(r,l)\mapsto v],(),k\rangle\,;\,\Write(l\at r,v)}
\end{mathpar}
\end{formalbox}

\paragraph{Effect-valued expression steps.}
Surface-effect terms are ordinary expressions whose values are computed effects.
Abstract summaries return region-level actions without touching the heap;
concrete summaries first evaluate their argument and then return a concrete
computed action. These rules are ordinary machine steps and therefore emit
$\Nil$.

\begin{formalbox}{Effect-Valued Expression Steps}
\footnotesize
\noindent\textbf{One machine step}\hfill $\boxed{s\longrightarrow s'\,;\,\eta}$
\begin{mathpar}
\inferrule*[lab=S-Abst-Read]{P(p)=r}
  {\langle H,E,P,\Rhat(p),k\rangle
   \longrightarrow
   \langle H,\Rs(r),k\rangle\,;\,\Nil}

\inferrule*[lab=S-Abst-Write]{P(p)=r}
  {\langle H,E,P,\What(p),k\rangle
   \longrightarrow
   \langle H,\Ws(r),k\rangle\,;\,\Nil}

\inferrule*[lab=S-Abst-Alloc]{P(p)=r}
  {\langle H,E,P,\Ahat(p),k\rangle
   \longrightarrow
   \langle H,\As(r),k\rangle\,;\,\Nil}

\inferrule*[lab=S-Conc-Read]{ }
  {\langle H,l\at r,K^{read}(k)\rangle
   \longrightarrow
   \langle H,\Read(l\at r),k\rangle\,;\,\Nil}

\inferrule*[lab=S-Conc-Write]{ }
  {\langle H,l\at r,K^{write}(k)\rangle
   \longrightarrow
   \langle H,\Write(l\at r),k\rangle\,;\,\Nil}

\inferrule*[lab=S-Concat]{ }
  {\langle H,\Theta_2,K^{concat}(\Theta_1,k)\rangle
   \longrightarrow
   \langle H,\Theta_1\cup\Theta_2,k\rangle\,;\,\Nil}

\inferrule*[lab=S-Top]{ }
  {\langle H,E,P,\Top,k\rangle
   \longrightarrow
   \langle H,\top,k\rangle\,;\,\Nil}

\inferrule*[lab=S-Empty]{ }
  {\langle H,E,P,\Pure,k\rangle
   \longrightarrow
   \langle H,\emptyset,k\rangle\,;\,\Nil}
\end{mathpar}
\end{formalbox}

Effect-summary evaluation is proved read-only from typing: if a well-typed
summary execution emits a read-only trace, the final heap is the initial heap.
A summary may \emph{denote} an allocation through an abstract allocation action
without performing an allocation while the summary itself is evaluated.



\section{Static Semantics}
\label{sec:staticsem}

The typing rules for the ordinary surface-effect language are given in the
present section. A variable context $\Gamma$ is a total map that records free
variables and their types and a region context $\Delta$ is a set of region
variables in scope. The judgment
$\Gamma; \Delta \vdash e : \tau\,!\,\epsilon$ asserts that expression $e$ has
type $\tau$ under $\Gamma$ and $\Delta$ and the latent (static) effect
$\epsilon$. As declared in rule [T-ABST], the correctness between a
computational--effect pair is denoted by the binary relation ($\desc$). The
rule [T-CAPP] declares the return type of a computational application,
$\mathsf{capp}$, and the same reasoning is made for $\mathsf{eapp}$ (whose
return type is always $\mathsf{eff}$).

Because the mechanization uses logical equality when proving soundness
properties about surface-effect terms, types are represented in a locally nameless
style \citep{chargueraud2012}, so that quantified types are equivalent up to
$\alpha$-conversion. Bound variables are represented with de Bruijn indices
and free-variable names are left unchanged. This representation lets
substitution be stated directly by type-level $\beta$-reduction.

The mechanized predicate $\ReadOnly$ on static effects is generated by three
cases: the empty effect is read-only, a singleton region-read effect is
read-only, and unions of read-only effects are read-only. Ordinary
static-effect soundness is an auxiliary theorem about evaluation traces: if the
static erasure $\staticset{\phi}$ of a trace is contained in a read-only static
effect, then the trace itself contains only reads. This is separate from the
computed-effect precision bridge $\phi\approxx\Theta$, which deliberately keeps
concrete locations until they are erased for static comparison. The
corresponding dynamic-trace predicate contains only the empty trace, concrete
read actions, and sequential or parallel compositions of read-only traces;
allocation and write actions are excluded. Consequently, read-only summary
evaluation preserves the heap. This is a semantic property, not a cost claim:
the development does not prove that a read-only summary terminates or that
evaluating it has constant overhead.

\begin{corollary}[Empty static trace erasure implies the empty trace]
For all traces $\phi$, if $\staticset{\phi}$ is contained in the empty static
effect, then $\phi = \Nil$.
\end{corollary}

\begin{lemma}[Read-only traces of dynamic effects imply no side-effects]
Given that $\ReadOnly(\phi)$, a terminal structured execution
\[
\mathsf{init}(H,E,P,e)\longrightarrow^{*}\mathsf{done}(H',v)\,;\,\phi
\]
implies that $H \equiv H'$ up to key-value pairs.
\end{lemma}

\begin{lemma}[Read-only static effects imply read-only traces]
If $\staticset{\phi}\subseteq \epsilon$ and $\ReadOnly(\epsilon)$ holds for
static effects, then $\ReadOnly(\phi)$ holds for the dynamic trace.
\end{lemma}

\section{Correctness of Surface Effects}
\label{sec:correctness}

The previous section defined the typed bridge $\Theta\covers\phi$ between
computed effects and dynamic traces. This section explains how that bridge is
obtained from source-level verification. The judgment
\[
  \Gamma;\Delta;P \vdash e \desc \thetahat
\]
relates a computation $e$ to a surface-effect term $\thetahat$. If $e$ executes
with trace $\phi$ and $\thetahat$ evaluates to a computed effect $\Theta$, then
surface-effect correctness requires
\[
  \Theta\covers\phi,
\]
equivalently $\phi\approxx\Theta$.  The rules defining this coverage relation
are part of the general effect vocabulary in Section~\ref{sec:effects}; this
section proves that verified surface-effect terms establish that relation.

The judgment is intentionally syntax directed but incomplete: there are
infinitely many sound summaries for a given expression, and some expressions can
always be over-described by `$\Top$'.  The theorem below therefore states a
soundness property for derivable descriptions, not a completeness result for all
possible summaries.


The correctness theorem below follows the clean direction of the design.
A surface-effect term evaluates to a computed effect, and the computation emits
a trace.  The theorem proves that the already-computed effect covers the
computation trace:
\[
\widehat{\Theta}\longrightarrow^{*}\Theta
\qquad\text{and}\qquad
\phi\approxx\Theta.
\]
The proof keeps traces and summaries as separate objects.  If evaluating a
function position, an argument position, or another subexpression may perform
memory actions, those actions must be accounted for by the surface-effect term
and by the source-level correctness judgment itself; they are not recovered
after the fact.

\begin{formalbox}{Back-Triangle Rules}
\noindent\textbf{Back-triangle rules ($\Gamma; \Delta; P \vdash e \desc \thetahat$).}

\noindent\emph{In words.}  The judgment says that the source term $\thetahat$ is an executable summary for the computation $e$ under the current region map $P$.
The judgment is indexed by the run-time region map $P$. For a static effect
$\epsilon$, write $\epsilon[P]$ for the effect obtained by substituting the
region values in $P$. The following rules give the paper-facing inductive
presentation. In particular, pure constants, booleans, variables, abstractions,
and region abstractions are described by $\Pure$, rather than by an arbitrary
summary.

\begin{mathpar}
\inferrule*[lab=Eff-Const]
  {\Gamma;\Delta \vdash \mathsf{const}\; n : \mathsf{nat}\,!\,\emptyset}
  {\Gamma;\Delta;P \vdash \mathsf{const}\; n \desc \Pure}

\inferrule*[lab=Eff-Bool]
  {\Gamma;\Delta \vdash \mathsf{bool}\; b : \mathsf{bool}\,!\,\emptyset}
  {\Gamma;\Delta;P \vdash \mathsf{bool}\; b \desc \Pure}

\inferrule*[lab=Eff-Var]
  {\Gamma;\Delta \vdash \mathsf{var}\; x : \tau\,!\,\emptyset}
  {\Gamma;\Delta;P \vdash \mathsf{var}\; x \desc \Pure}

\inferrule*[lab=Eff-Abst]{ }
  {\Gamma;\Delta;P \vdash \mu f. \lambda x. \{e_b \mid e_e\} \desc \Pure}
\end{mathpar}

\begin{mathpar}
\inferrule*[lab=Eff-Rgn-Abst]{ }
  {\Gamma;\Delta;P \vdash \lambda \rho.\,e \desc \Pure}

\inferrule*[lab=Eff-App]
  {\Gamma;\Delta \vdash \capp e_f\,e_a : \tau_b\,!\,\epsilon_b \\
   \Gamma;\Delta \vdash \eapp e_f\,e_a : \mathsf{eff}\,!\,\epsilon_e \\
   \ReadOnly(\epsilon_e[P]) \\
   \Gamma;\Delta \vdash e_f : \tau_f\,!\,\epsilon_f \\
   \Gamma;\Delta \vdash e_a : \tau_a\,!\,\epsilon_a \\
   \Gamma;\Delta;P \vdash e_f \desc (\eapp e_f\,e_a) \\
   \Gamma;\Delta;P \vdash e_a \desc (\eapp e_f\,e_a) \\
   \ReadOnly(\epsilon_f[P]) \\
   \ReadOnly(\epsilon_a[P])}
  {\Gamma;\Delta;P \vdash \capp e_f\,e_a \desc (\eapp e_f\,e_a)}

\inferrule*[lab=Eff-Rgn-App]
  {\Gamma;\Delta \vdash e_r : \tau\,!\,\epsilon \\
   \Gamma;\Delta \vdash \Pure : \mathsf{eff}\,!\,\emptyset \\
   \Gamma;\Delta;P \vdash e_r \desc \Pure}
  {\Gamma;\Delta;P \vdash e_r\,[r] \desc \Pure}

\inferrule*[lab=Eff-Cond]
  {\Gamma;\Delta \vdash e_c : \tau_c\,!\,\epsilon_c \\
   \Gamma;\Delta \vdash e_t : \tau_t\,!\,\epsilon_t \\
   \Gamma;\Delta \vdash e_f : \tau_f\,!\,\epsilon_f \\
   \ReadOnly(\epsilon_c[P]) \\
   \Gamma;\Delta;P \vdash e_c \desc \Pure \\
   \Gamma;\Delta;P \vdash e_t \desc e_1 \\
   \Gamma;\Delta;P \vdash e_f \desc e_2}
  {\Gamma;\Delta;P \vdash \cond e_c\,e_t\,e_f \desc
   \cond e_c\,e_1\,e_2}

\inferrule*[lab=Eff-Alloc]
  {\Gamma;\Delta \vdash e : \tau\,!\,\epsilon \\
   \Gamma;\Delta;P \vdash e \desc e_1}
  {\Gamma;\Delta;P \vdash \alloc r\;e \desc e_1 \oplus \Ahat(r)}

\inferrule*[lab=Eff-Get-Abs]
  {\Gamma;\Delta \vdash e : \tau\,!\,\epsilon \\
   \Gamma;\Delta;P \vdash e \desc e_1}
  {\Gamma;\Delta;P \vdash \getop p\;e \desc e_1 \oplus \Rhat(p)}
\end{mathpar}

\begin{mathpar}
\inferrule*[lab=Eff-Get-Conc]
  {\Gamma;\Delta \vdash e : \tau\,!\,\epsilon \\
   \Gamma;\Delta;P \vdash e \desc \Pure}
  {\Gamma;\Delta;P \vdash \getop r\;e \desc \widehat{\Read(e)}}

\inferrule*[lab=Eff-Set-Abs]
  {\Gamma;\Delta;P \vdash e_r \desc e_1 \\
   \Gamma;\Delta;P \vdash e_v \desc e_2 \\
   \Gamma;\Delta \vdash e_r : \tau_r\,!\,\epsilon_r \\
   \ReadOnly(\epsilon_r[P])}
  {\Gamma;\Delta;P \vdash \setop p\;e_r\,e_v \desc
   e_1 \oplus (e_2 \oplus \What(p))}

\inferrule*[lab=Eff-Set-Conc]
  {\Gamma;\Delta;P \vdash e_r \desc e_1 \\
   \Gamma;\Delta;P \vdash e_v \desc e_2 \\
   \Gamma;\Delta \vdash e_r : \tau_r\,!\,\epsilon_r \\
   \ReadOnly(\epsilon_r[P])}
  {\Gamma;\Delta;P \vdash \setop r\;e_r\,e_v \desc
   e_1 \oplus (e_2 \oplus \Writehat(e_r))}

\inferrule*[lab=Eff-Arith]
  {\Gamma;\Delta \vdash e_1 : \tau_1\,!\,\epsilon_1 \\
   \Gamma;\Delta \vdash e_2 : \tau_2\,!\,\epsilon_2 \\
   \ReadOnly(\epsilon_1[P]) \\
   \Gamma;\Delta;P \vdash e_1 \desc \thetahat_1 \\
   \Gamma;\Delta;P \vdash e_2 \desc \thetahat_2}
  {\Gamma;\Delta;P \vdash e_1\,\mathbin{\odot}\,e_2
   \desc \thetahat_1 \oplus \thetahat_2}
\end{mathpar}

\begin{mathpar}
\inferrule*[lab=Eff-Top]{ }
  {\Gamma;\Delta;P \vdash e \desc \Top}
\end{mathpar}

Here $\odot \in \{+,-,\times,=\}$ abbreviates the four corresponding arithmetic
and equality cases. The concrete read and write rules above use a concrete
region value $r$.
\end{formalbox}

\noindent\emph{In words.} A verified surface-effect term evaluates to a computed effect that covers the entire trace of the paired computation.

\begin{theorem}[Correctness of surface-effect terms]
Assume well-formed region, typing, heap-typing, and environment-typing
judgments for $H$, $E$, $P$, and $e_b$. If the computation and summary are
related by
$\Gamma;\Delta;P \vdash e_b \desc e_e$, and
\[
\mathsf{init}(H,E,P,e_b)\longrightarrow^{*}\mathsf{done}(H_b,v_b)\,;\,\phi_b
\qquad\text{and}\qquad
\mathsf{init}(H,E,P,e_e)\longrightarrow^{*}
  \mathsf{done}(H_e,\mathsf{eff}(\Theta_e))\,;\,\phi_e,
\]
with $\ReadOnly(\phi_e)$, then
\[
\Theta_e\covers\phi_b
\qquad\text{equivalently}\qquad
\phi_b \approxx \Theta_e.
\]
\end{theorem}

The conclusion of the correctness theorem is exactly the right-hand typed
bridge. In particular, it does \emph{not} conclude that the computation trace
$\phi_b$ is read-only, nor does it prove the left-hand bridge
$\epsilon\covers\Theta_e$. The small-step results in Section~\ref{sec:smallstep}
reuse this approximation story while stating progress, preservation, terminal
determinism, and terminal correctness directly over the continuation machine.

\begin{corollary}[Read-only, heap-preserving summary evaluation]
For a well-typed computation/effect application pair, static effect soundness
and the typing premises of the correctness relation imply
$\ReadOnly(\phi_e)$. Therefore summary evaluation preserves the initial heap:
$H \equiv H_e$.
\end{corollary}

The proof separates two facts: the computation trace is shown to satisfy
$\phi_b \approxx \Theta_e$, while the trace produced by evaluating the effect
application is separately shown to be read-only. Thus a computed action such as
$\As(r)$ may \emph{describe} a future allocation without the evaluation of the
summary itself allocating memory.

\paragraph{Trace/summary separation.}
The previous theorem is the only coverage principle needed for surface-effect
correctness: a verified surface-effect term must cover the whole trace of its
paired computation.  The dynamic-action views $\dynset{\phi}$ and
$\staticset{\phi}$ are used only to state erasure and coverage facts about
already produced traces.

\noindent\emph{In words.} If two traces are covered by disjoint and
nonconflicting computed effects, then their parallel trace composition is
deterministic.

\begin{theorem}[Trace-composition soundness]
Let $\phi_1 \approxx \Theta_1$ and $\phi_2 \approxx \Theta_2$. If
\[
\Disjoint(\Theta_1,\Theta_2)
\qquad\text{and}\qquad
\neg(\Theta_1 \conflict \Theta_2),
\]
and $\DetTrace(\phi_1)$ and $\DetTrace(\phi_2)$, then
\[
\DetTrace(\phi_1 \,\|\, \phi_2).
\]
\end{theorem}

The proof first shows that a dynamic read--write or write--write
conflict would induce a corresponding computed conflict, contradicting the
summary-level premise. It separately transports computed-action disjointness to
every cross-trace pair of dynamic actions. The result is the deterministic
parallel trace judgment above.

Combining the preceding results gives the trace-level proof dependency:
\[
\begin{array}{c}
\text{verified computation/summary pair}
\\[1mm]\Longrightarrow\\[-1mm]
\phi_i \approxx \Theta_i
\\[1mm]\Longrightarrow\;\Disjoint(\Theta_1,\Theta_2)\land
  \neg(\Theta_1\conflict\Theta_2)\\[-1mm]
\DetTrace(\phi_1\,\|\,\phi_2)
\\[1mm]\Longrightarrow\\[-1mm]
\text{equivalent final heaps under terminating trace replay.}
\end{array}
\]

\paragraph{Runtime definitions.}
\[
\Sigma \triangleq \mathit{Addr} \rightharpoonup \mathit{Type}
\qquad \text{(heap type assignment)}
\]

\begin{mathpar}
\inferrule*[lab=T-Heap]
  {\forall l\at r \in H,\; \Sigma \vdash H(l\at r) : \Sigma(l\at r)}
  {\star \vdash H : \Sigma}

\inferrule*[lab=T-Heap-Empty]{ }
  {\Sigma \vdash \emptyset : \emptyset}

\inferrule*[lab=T-Heap-Ext]
  {\Sigma \vdash E : \Gamma \\ \Sigma \vdash v : \tau}
  {\Sigma \vdash E, [x \mapsto v] : \Gamma, v : \tau}
\end{mathpar}

As previously mentioned, the static relation ($\desc$) is intentionally
incomplete: many sound summaries may describe the same computation, while some
computations may be related only to a coarse summary by the supplied inductive
rules. Concrete computed actions can refine a region-level distinction and may
therefore satisfy computed-effect disjointness even when a region-only
comparison would be inconclusive. `$\Top$', by contrast, is a sound
approximation of every trace but cannot provide useful disjointness evidence.
The current calculus contains no separate static recovery rule for `$\Top$';
its role in the mechanized semantics is to conservatively lose precision.

For example, the main purpose of designing a type system where the functional
type contains a return type for a computational expression and a return type
for an effect expression is to be able to describe any computational
application ($\capp e_f\, e_a$) by its corresponding effect application
($\eapp e_f\, e_a$). The interesting aspect of assigning the ``extended''
functional type to the expression $e_f$ is that it may carry a non-empty
latent effect for the computational term which can be soundly described by the
evaluation of the corresponding ``surface'' effect term. The returned computed
summary of $\eapp e_f\,e_a$ is produced by the function body's summary
expression.  If evaluating $e_f$ or $e_a$ may itself perform memory actions,
those actions must be covered by the surface-effect term assigned to the whole
computation by the correctness judgment.  They are not converted into computed
summary components after the fact.

Another example is the expression $e_b$ whose reference type is parametrized
by a region $p$. Two different effect summaries are proven sound: the first is
as much precise as the static effect would be (rule [EFF-GET]) and the second
is potentially ``much'' more precise than the static effect (rule
[EFF-GET-CONC]). The type of the argument of $\mathsf{get}$ is the reference
type parametrized by a region variable or a region constant. Given the region
$p$, the rule [EFF-GET] applies by chaining the effect summary of the argument
with the static surface term $\Rhat(p)$. When evaluated, this term yields a
value that is equivalent to a de facto static effect. However, we can
precisely describe a $\mathsf{get}$ operation using the surface-effect term
$\widehat{\Read(e)}$, which directly takes the argument as input in order to evaluate
to a concrete memory address inside some allocated region.

The precision of surface-effect terms can be evaluated by analyzing their
dynamic semantics. Consider, for example, the rules [S-Abst-Read] and
[S-Conc-Write] in Section~\ref{sec:opsem}. The first corresponds to the
evaluation of the ``static version'' of the surface-effect term for the
$\mathsf{get}$ operation and the second corresponds to the evaluation of the
``concrete version'' of the surface-effect term for the $\mathsf{set}$ operation.
The second rule shows that a concrete address, $l\at r$, is being evaluated.
Now consider the rules for conflict decision presented in
Section~\ref{sec:effects}. When concrete allocated addresses are known, it is
possible to decide that the effects $\Read$ and $\Write$ do not conflict if
the addresses are different even though they belong to the same region,
allowing more precision (hence the motivating example presented in
Fig.~\ref{fig:example}). However, if only the region is known, the ``static
version'' of the summary, $\Rs$ and $\Ws$, does not allow the conflict decision
at run-time to be more precise than the static analysis.

\section{Small-Step Runtime Safety}
\label{sec:smallstep}

The continuation machine supports safety statements over finite executions,
not only over completed runs. The ordinary small-step safety invariant is
indexed by a store typing
$\Sigma$. Along a finite trace, allocation may extend $\Sigma$; preservation
therefore produces a final store typing $\Sigma'$ with
$\Sigma \subseteq \Sigma'$. We write
$s\longrightarrow^{*} s'\,;\,\phi$ for a finite small-step execution whose
transition fragments accumulate to $\phi$, $\mathsf{init}(H,E,P,e)$ for the initial machine state,
and $\mathsf{done}(H,v)$ for a terminal state. The judgment
$\Sigma \vdash_{\mathsf{st}} s : \tau$ abbreviates the mechanized typed-state
invariant, and $\Sigma \vdash_{\mathsf{tr}} \phi$ states that the trace
$\phi$ is well typed with respect to $\Sigma$.
The predicate $\mathsf{NotStuck}(s)$ means that $s$ is terminal or has a next
small-step transition. The predicate $\mathsf{Shape}_{\Sigma}(v,\tau)$ records
that the run-time form of $v$ matches type $\tau$ under store typing $\Sigma$.

\noindent\emph{In words.} Well-typed ordinary machine states do not get stuck,
and ordinary machine steps preserve the typed-state invariant.

\begin{theorem}[Progress]
If $\Sigma\vdash_{\mathsf{st}} s : \tau$, then either $s$ is terminal or there
exist $s'$ and $\eta$ such that
\[
s\longrightarrow s'\,;\,\eta.
\]
\end{theorem}

\begin{theorem}[One-step preservation]
If
\[
\Sigma\vdash_{\mathsf{st}} s : \tau
\qquad\text{and}\qquad
s\longrightarrow s'\,;\,\eta,
\]
then $\Sigma\vdash_{\mathsf{st}} s' : \tau$.
\end{theorem}

\begin{theorem}[Terminal trace determinism]
For any ordinary initial state $s$, if
\[
s\longrightarrow^{*}\mathsf{done}(H_1,v_1)\,;\,\phi_1
\qquad\text{and}\qquad
s\longrightarrow^{*}\mathsf{done}(H_2,v_2)\,;\,\phi_2,
\]
then
\[
\phi_1=\phi_2
\qquad\land\qquad
H_1=H_2
\qquad\land\qquad
v_1=v_2.
\]
\end{theorem}

\begin{theorem}[Terminal surface-effect correctness]
Assume the counted-run correctness invariant holds for every natural bound.
If
\[
\Gamma;\Delta\vdash e\desc \thetahat,
\qquad
\mathsf{init}(H,E,P,e)
  \longrightarrow^{*}
  \mathsf{done}(H',v)\,;\,\phi,
\]
and the summary expression terminates with computed effect $\Theta$,
\[
\mathsf{init}(H,E,P,\thetahat)
  \longrightarrow^{*}
  \mathsf{done}(H_s,\mathsf{eff}(\Theta))\,;\,\phi_s,
\]
then the computation trace is covered by the computed effect:
\[
\phi\approxx\Theta.
\]
\end{theorem}

The structured-trace version follows from the raw trace statement by erasing
the structured view to its underlying machine trace. The premise about the
counted-run correctness invariant is the mechanized induction closure used by
the ordinary small-step proof; it is not an operational assumption and it is
not evidence extracted from a trace.

\section{Discussion and Related Work}
\label{sec:related}

The closest static comparator is Deterministic Parallel Java (DPJ)
\citep{bocchino2009}. DPJ uses a modular compile-time type-and-effect system to
guarantee deterministic semantics and expresses effects over regions. Surface
effects retain static verification but move a value-dependent access summary
into an executable source-level term. From the run-time side, classic
inspector--executor work analyzes irregular dependences before the executor
runs conflict-free work in parallel \citep{arenaz2004}. The distinctive point
of the present formulation is therefore not region effects or pre-execution
access prediction in isolation, but the integration of a programmer-written,
typed summary term with an approximation theorem and a trace-composition
proof.

This positioning also exposes a limitation. Expressive region/effect
annotations can be difficult to write manually, and prior work has investigated
inference precisely to reduce that burden \citep{tzannes2019}. The present
calculus assumes programmer-written surface-effect summaries and does not give
an inference or synthesis algorithm. Annotation effort and the frequency with
which summaries collapse to `$\Top$' are empirical questions left open here.

Hybrid verification platforms are also close in spirit
\citep{kawaguchi2012,flanagan2006}.  They use expressive specifications, such
as refinement types, to combine static reasoning with value-sensitive
properties.  Liquid Types support deterministic fine-grained multithreading by
expressing effects in first-order formulas over a decidable arithmetic theory.
At a higher-level object-oriented setting, HOOL supports expressive object
specifications, including subtyping, refinements, and object invariants.  As in
our calculus, these specifications are pure: evaluating them does not update
mutable data.

The design is hybrid in a precise sense. Verification of the
computation/summary relation is static, while the summary value is dynamic:
the summary term is evaluated with run-time values and its computed effect can
be checked for disjointness and conflict absence by later clients. Summary
evaluation is still run-time computation and may itself be input dependent.
The continuation-machine metatheory proves progress, preservation, terminal
determinism, and terminal surface-effect correctness for ordinary executions.
It does not prove termination of summary terms, bound their cost, or compare
that cost with dynamic validation mechanisms. A future adequacy theorem could
relate terminal machine runs to the archived terminating evaluator, but that
bridge is not used as a paper theorem here.

Recent work on hybrid type and effect systems applied to concurrent
programming is found in \citet{long2013}. The approach defines special static
effects, designated open effects, that prevent object-oriented programs from
being rejected when no static upper bound in the form of effect specification
is available. More specifically, open effects are assumed not to conflict with
any other effects and are used when the dynamic type of receiver methods
cannot be known statically due to dynamically dispatched method invocation.
Although open effects are somewhat analogous to the surface-effects approach,
in the sense that disjointness information may be unknown until runtime,
there is a subtle difference: while open effects are inductively defined and designed for
optimistic parallelism, where concretized effects soundly approximate runtime
effects, surface-effect terms are syntactically defined in terms of unevaluated
expressions and their concretization derives from runtime heap allocation.
The correctness of surface-effect terms is always statically known, and
the purpose of dynamic checking is to find further possibilities for parallelism.
Moreover, surface-effect terms can be used in those scenarios where conflict
decision based on static effects would reject parallel programs. Since more
precise effect specifications, yet statically verified, are provided directly
by the programmer, surface-effect analysis can override static-effect
analysis.

Irregular data structures such as trees, queues, and graphs motivate the use of
run-time information because region-level analyses may merge accesses that are
concrete-address disjoint. The Galois approach \citep{kulkarni2007} is a useful
contrast because it uses run-time conflict detection in an optimistic execution
model. The surface-effects approach instead evaluates a summary before the parallel
computation and therefore do not require the calculus to roll back a
computation after a summary-level conflict is found. This is a semantic
difference, not an established performance comparison.

The present work supplies neither a cost semantics nor an implementation study.
Consequently, we make no claim that summary evaluation is constant in the size
of the summary, that pre-checking has lower total overhead than access logging,
or that useful parallelism is recovered frequently in practice. A future
evaluation should measure at least summary-evaluation time, conflict-test time,
false-positive rate, the frequency of `$\Top$', recovered parallelism relative
to region-only static effects, and programmer annotation effort. These
measurements are necessary to determine whether the formal precision mechanism
translates into a practical advantage.

\section{Conclusions}
\label{sec:conclusions}

We have presented a region-polymorphic calculus in which a computation is
paired with a statically verified, executable effect-summary term. Evaluating
the summary before the computation can expose concrete region/location actions
that a region-only static effect does not distinguish. The mechanized
correctness theorem proves trace approximation, and the trace-composition
lemma makes the bridge explicit: disjoint and non-conflicting computed
summaries of deterministic component traces yield a deterministic parallel
trace. The trace semantics then establishes equivalent terminating replay
heaps. The continuation-machine theorem stack adds ordinary runtime safety:
well-typed states make progress, machine steps preserve typing, and two
terminal executions from the same initial state have the same trace, heap, and
value.

These results are formal, not empirical performance results. The calculus does
not include summary inference, a termination argument for summary evaluation, a
cost semantics, an implementation evaluation, or a separate adequacy theorem for
terminal continuation-machine runs.
Establishing the annotation burden and the run-time trade-off between summary
precision, pre-checking cost, and recovered parallelism is therefore left to
future work.

\bibliographystyle{plainnat}
\bibliography{references}

\appendix

\section{Deterministic Parallelism}
\label{app:determinism}

\begin{definition}[Heap equivalence]
The mechanization represents a heap as a finite map from keys $(r,l)$ to
values. We write $H_1 \equiv H_2$ when the two finite maps are extensionally
equal, equivalently when every key has the same lookup result in both heaps:
\[
H_1 \equiv H_2
\quad\Longleftrightarrow\quad
\forall k.\; H_1(k)=H_2(k).
\]
This is equality of key--value maps, not an isomorphism modulo renaming of
locations. The relation is reflexive, symmetric, and transitive.
\end{definition}

\begin{definition}[Deterministic trace]
The predicate $\DetTrace(\phi)$ is inductively generated by the empty trace,
single dynamic actions, sequential composition of deterministic traces, and a
parallel constructor. The latter requires deterministic component traces and
both
\[
\neg\mathsf{ConflictTraces}(\phi_1,\phi_2)
\qquad\text{and}\qquad
\mathsf{DisjointTraces}(\phi_1,\phi_2).
\]
Thus absence of an explicitly represented read--write or write--write conflict
is not the only premise: every pair of dynamic actions drawn from opposite
branches must also satisfy dynamic-action disjointness.
\end{definition}

\begin{definition}[Trace replay]
The relation $\Rightarrow$ executes one dynamic action of a trace against a
heap. Allocation and write actions update the heap, reads require the recorded
value to be present, and sequential/parallel trace nodes provide the evaluation
contexts. We write $\Rightarrow^{*}$ for its reflexive transitive closure.
A terminating replay has the form
$(\phi,H)\Rightarrow^{*}(\Nil,H')$.
\end{definition}

\begin{lemma}[Read-only trace properties]
A read-only trace preserves the heap during replay. Two read-only component
traces form a deterministic parallel trace. These results follow from the
inductive definition of $\ReadOnly$ on traces, which contains only the empty
trace, read actions, and sequential/parallel compositions of read-only traces.
\end{lemma}

\begin{lemma}[Scheduler soundness]
For $i\in\{1,2\}$, let $\phi_i \approxx \Theta_i$ and
$\DetTrace(\phi_i)$. If
\[
\Disjoint(\Theta_1,\Theta_2)
\quad\land\quad
\neg(\Theta_1 \conflict \Theta_2),
\]
then $\DetTrace(\phi_1\,\|\,\phi_2)$.
\end{lemma}

The proof uses the soundness relation $\approxx$ in both directions required by
the parallel-trace constructor: a dynamic conflict would imply a computed
conflict, and computed-action disjointness implies disjointness of every
represented pair of dynamic actions.

\begin{theorem}[Confluence of deterministic trace replay]
Let $H_a \equiv H_b$. If a deterministic trace $\phi$ is replayed from the two
equivalent heaps to intermediate configurations
$(\phi_1,H_1)$ and $(\phi_2,H_2)$, then there exist a common trace $\phi_3$ and
heaps $H_3 \equiv H_4$ such that
\[
(\phi_1,H_1)\Rightarrow^{*}(\phi_3,H_3)
\qquad\text{and}\qquad
(\phi_2,H_2)\Rightarrow^{*}(\phi_3,H_4).
\]
In particular, if both replays terminate at $\Nil$, their final heaps are
equivalent.
\end{theorem}

The result does not claim uniqueness of an intermediate trace independently of
the stated confluence property.

\begin{theorem}[Heap uniqueness for safe parallel summaries]
Let $\phi_1 \approxx \Theta_1$ and $\phi_2 \approxx \Theta_2$, with deterministic
component traces and
\[
\Disjoint(\Theta_1,\Theta_2)
\quad\land\quad
\neg(\Theta_1 \conflict \Theta_2).
\]
Let $H_a \equiv H_b$. If
\[
(\phi_1\,\|\,\phi_2,H_a)\Rightarrow^{*}(\Nil,H_1)
\quad\text{and}\quad
(\phi_1\,\|\,\phi_2,H_b)\Rightarrow^{*}(\Nil,H_2),
\]
then $H_1 \equiv H_2$.
\end{theorem}

The proof applies trace-composition soundness to establish determinism of the parallel
trace and then uses the terminating replay theorem.

\end{document}
